\documentclass[10pt,a4paper,twocolumn]{article}
\usepackage[utf8]{inputenc}
\usepackage[T1]{fontenc}
\usepackage{geometry}
\usepackage{amsmath,amssymb}
\usepackage{booktabs,array,tabularx}
\usepackage{enumitem}
\usepackage{listings}
\usepackage{xcolor}
\usepackage{caption}
\geometry{left=0.9cm,right=0.9cm,top=1.0cm,bottom=1.0cm,columnsep=0.55cm}
\setlength{\parindent}{0pt}
\setlength{\parskip}{2pt}
\setlist{nosep,leftmargin=*}
\captionsetup{font=scriptsize,labelfont=bf}
\lstset{basicstyle=\ttfamily\scriptsize,breaklines=true,columns=fullflexible,frame=single,showstringspaces=false,keepspaces=true}
\newcommand{\jawab}{\textbf{Jawab: }}
\newcommand{\soal}{\textbf{Soal: }}
\newcommand{\yes}{\textsc{Yes}}
\newcommand{\no}{\textsc{No}}
\sloppy

\title{Pembahasan Ringkas Latihan Soal UAS Machine Learning}
\author{Bahan Belajar - Format 2 Kolom}
\date{Semester II 2024/2025}

\begin{document}
\maketitle

\section*{1. Data Preparation dan Evaluation}

\subsection*{1a. Train-Test Split}
\soal Dari kode pada Gambar 1, jika total data adalah 500 dan \texttt{train\_size = 0.7}, tuliskan output program.

\textbf{Deskripsi Gambar 1:} potongan kode Python memakai \texttt{train\_test\_split(X,y, train\_size=0.7)} lalu mencetak jumlah data latih dan data uji.

\begin{lstlisting}[language=Python]
from sklearn.model_selection import train_test_split
X_train, X_test, y_train, y_test = train_test_split(X, y, train_size=0.7)
print("Banyak data latih setelah dilakukan Train-Test Split: ", len(X_train))
print("Banyak data uji setelah dilakukan Train-Test Split: ", len(X_test))
\end{lstlisting}

\jawab
\[
|X_{train}|=0.7\times 500=350, \qquad |X_{test}|=500-350=150.
\]
Output:
\begin{lstlisting}
Banyak data latih setelah dilakukan Train-Test Split: 350
Banyak data uji setelah dilakukan Train-Test Split: 150
\end{lstlisting}

\subsection*{1b. K-Fold Cross Validation}
\soal Jelaskan \textit{k-fold cross validation}. Kapan digunakan?

\jawab \textit{k-fold cross validation} membagi data menjadi $k$ bagian/fold. Model dilatih pada $k-1$ fold dan diuji pada 1 fold secara bergiliran, lalu skor dirata-ratakan. Dipakai ketika data terbatas, ingin evaluasi model lebih stabil, dan ingin membandingkan model/hyperparameter tanpa terlalu bergantung pada satu pembagian train-test.

\subsection*{1c. Confusion Matrix dan Metrik}
\soal Dari Tabel 1 berikut, buat confusion matrix lalu hitung accuracy, sensitivity/recall, precision, dan F1-measure. Kelas positif adalah \yes.

\begin{center}
\scriptsize
\begin{tabular}{@{}ccc@{}}
\toprule
No & Target & Prediksi \\
\midrule
1 & \no & \no \\
2 & \no & \yes \\
3 & \yes & \yes \\
4 & \yes & \no \\
5 & \yes & \yes \\
6 & \no & \no \\
7 & \no & \no \\
8 & \yes & \no \\
9 & \yes & \yes \\
10 & \yes & \no \\
11 & \yes & \no \\
12 & \yes & \yes \\
\bottomrule
\end{tabular}
\end{center}

\jawab
\begin{center}
\scriptsize
\begin{tabular}{@{}lcc@{}}
\toprule
 & Pred. \yes & Pred. \no \\
\midrule
Aktual \yes & TP = 4 & FN = 4 \\
Aktual \no & FP = 1 & TN = 3 \\
\bottomrule
\end{tabular}
\end{center}

\[
\text{Accuracy}=\frac{TP+TN}{12}=\frac{7}{12}=58.33\%.
\]
\[
\text{Recall}=\frac{TP}{TP+FN}=\frac{4}{8}=50\%.
\]
\[
\text{Precision}=\frac{TP}{TP+FP}=\frac{4}{5}=80\%.
\]
\[
F1=\frac{2PR}{P+R}=\frac{2(0.8)(0.5)}{0.8+0.5}=61.54\%.
\]

\section*{2. K-Nearest Neighbor (K-NN)}

\subsection*{2a. K-NN Cocok untuk Dataset Besar karena Prediksi Cepat}
\soal Tentukan benar/salah dan berikan alasan.

\jawab \textbf{Salah.} K-NN adalah \textit{lazy learner}; saat prediksi, model menghitung jarak data baru ke banyak/semua data latih. Jika dataset besar, prediksi cenderung lambat dan butuh memori besar.

\subsection*{2b. K-NN Harus Memakai Euclidean Distance}
\soal Tentukan benar/salah dan berikan alasan.

\jawab \textbf{Salah.} Euclidean distance hanya salah satu pilihan. K-NN bisa memakai Manhattan, Minkowski, cosine similarity, Hamming distance, atau metrik lain sesuai tipe data.

\section*{3. Naive Bayes: Play Tennis}

\subsection*{3a. Probability Table dan Prediksi D15}
\soal Diberikan data Play Tennis. Buat tabel probabilitas, lalu tentukan apakah pada D15 disarankan bermain tenis untuk kondisi: Outlook = Rain, Temperature = Mild, Humidity = Normal, Wind = Weak. Soal menyebut D15 dan D16, tetapi atribut cuaca yang tertulis hanya D15.

\begin{center}
\scriptsize
\begin{tabular}{@{}clllll@{}}
\toprule
Day & Outlook & Temp. & Humid. & Wind & Play \\
\midrule
D1 & Sunny & Hot & High & Weak & No \\
D2 & Sunny & Hot & High & Strong & No \\
D3 & Overcast & Hot & High & Weak & Yes \\
D4 & Rain & Mild & High & Weak & Yes \\
D5 & Rain & Cool & Normal & Weak & Yes \\
D6 & Rain & Cool & Normal & Strong & No \\
D7 & Overcast & Cool & Normal & Strong & Yes \\
D8 & Sunny & Mild & High & Weak & No \\
D9 & Sunny & Cool & Normal & Weak & Yes \\
D10 & Rain & Mild & Normal & Weak & Yes \\
D11 & Sunny & Mild & Normal & Strong & Yes \\
D12 & Overcast & Mild & High & Strong & Yes \\
D13 & Overcast & Hot & Normal & Weak & Yes \\
D14 & Rain & Mild & High & Strong & No \\
\bottomrule
\end{tabular}
\end{center}

\jawab Jumlah kelas: Play = \yes{} ada 9, Play = \no{} ada 5.

\begin{center}
\scriptsize
\begin{tabular}{@{}lcc@{}}
\toprule
Probabilitas & Play=\yes & Play=\no \\
\midrule
$P(C)$ & $9/14$ & $5/14$ \\
Outlook=Sunny & $2/9$ & $3/5$ \\
Outlook=Overcast & $4/9$ & $0/5$ \\
Outlook=Rain & $3/9$ & $2/5$ \\
Temp=Hot & $2/9$ & $2/5$ \\
Temp=Mild & $4/9$ & $2/5$ \\
Temp=Cool & $3/9$ & $1/5$ \\
Humidity=High & $3/9$ & $4/5$ \\
Humidity=Normal & $6/9$ & $1/5$ \\
Wind=Weak & $6/9$ & $2/5$ \\
Wind=Strong & $3/9$ & $3/5$ \\
\bottomrule
\end{tabular}
\end{center}

Untuk D15 = Rain, Mild, Normal, Weak:
\[
S_{Yes}=\frac{9}{14}\cdot\frac{3}{9}\cdot\frac{4}{9}\cdot\frac{6}{9}\cdot\frac{6}{9}=0.04233.
\]
\[
S_{No}=\frac{5}{14}\cdot\frac{2}{5}\cdot\frac{2}{5}\cdot\frac{1}{5}\cdot\frac{2}{5}=0.00457.
\]
Karena $S_{Yes}>S_{No}$, maka prediksi D15 adalah \textbf{Play Tennis = \yes}.

\subsection*{3b. Naive Bayes untuk Data Numerik}
\soal Naive Bayes awalnya untuk data kategorikal, tetapi bisa untuk data numerik. Jelaskan caranya.

\jawab Untuk fitur numerik, gunakan \textbf{Gaussian Naive Bayes}. Untuk setiap kelas dihitung mean $\mu$ dan varians $\sigma^2$, lalu likelihood fitur dihitung dengan distribusi normal:
\[
P(x\mid C)=\frac{1}{\sqrt{2\pi\sigma_C^2}}\exp\left(-\frac{(x-\mu_C)^2}{2\sigma_C^2}\right).
\]
Nilai likelihood tiap fitur dikalikan dengan prior kelas, lalu pilih kelas dengan posterior terbesar.

\section*{4. Ensemble Learning}

\subsection*{4a. Bagging, Boosting, dan Stacking}
\soal Jelaskan algoritma bagging, boosting, stacking, dan perbedaannya.

\jawab
\begin{itemize}
\item \textbf{Bagging}: membuat beberapa data bootstrap, melatih beberapa model secara paralel, lalu menggabungkan hasil dengan voting/average. Fokus mengurangi variance. Contoh: Random Forest.
\item \textbf{Boosting}: melatih model berurutan; model berikutnya fokus pada kesalahan model sebelumnya. Fokus mengurangi bias dan meningkatkan akurasi. Contoh: AdaBoost, Gradient Boosting.
\item \textbf{Stacking}: melatih beberapa model berbeda sebagai base model, lalu outputnya dipakai sebagai input model akhir/meta-learner.
\end{itemize}
Perbedaan utama: bagging paralel dan berbasis bootstrap; boosting sekuensial dan berbasis perbaikan error; stacking menggabungkan model berbeda memakai meta-model.

\subsection*{4b. Parameter AdaBoost}
\soal Dari Gambar 2, jelaskan arti \texttt{base\_classifier}, \texttt{n\_estimators}, dan \texttt{learning\_rate}.

\textbf{Deskripsi Gambar 2:} potongan kode membuat \texttt{AdaBoostClassifier(base\_classifier, n\_estimators=50, learning\_rate=1.0, random\_state=42)} lalu melakukan \texttt{fit(X\_train, y\_train)}.

\jawab
\begin{itemize}
\item \texttt{base\_classifier}: model dasar/weak learner yang dipakai AdaBoost, misalnya decision stump.
\item \texttt{n\_estimators=50}: jumlah weak learner yang dilatih secara berurutan.
\item \texttt{learning\_rate=1.0}: besar kontribusi tiap weak learner; semakin kecil biasanya lebih stabil tetapi butuh lebih banyak estimator.
\end{itemize}

\subsection*{4c. Bagging dengan 3 Model}
\soal Implementasikan bagging pada data berikut dengan 3 bootstrap/single model. Tulis prediksi akhir dan akurasi.

\begin{center}
\scriptsize
\begin{tabular}{@{}lccccccc@{}}
\toprule
$x$ & 2.4 & 3.1 & 3.7 & 4.4 & 5.2 & 5.8 & 6.4 \\
$y$ & 1 & 1 & 1 & -1 & -1 & -1 & -1 \\
\bottomrule
\end{tabular}
\end{center}

\jawab Karena sampel bootstrap acak tidak dicantumkan pada soal, dipakai asumsi base learner berupa decision stump. Data terpisah jelas pada batas tengah antara 3.7 dan 4.4:
\[
\theta=\frac{3.7+4.4}{2}=4.05.
\]
Aturan tiap model yang masuk akal:
\[
\hat{y}=\begin{cases}
1, & x<4.05,\\
-1, & x\ge 4.05.
\end{cases}
\]
Jika tiga model menghasilkan aturan yang sama, voting mayoritas memberi:
\[
\hat{y}=[1,1,1,-1,-1,-1,-1].
\]
Semua prediksi sama dengan label asli, sehingga:
\[
\text{Accuracy}=\frac{7}{7}\times 100\%=100\%.
\]
Catatan: jika bootstrap acak ditentukan berbeda oleh soal, prediksi model tunggal dapat berubah, tetapi prediksi akhir bagging tetap diambil dari voting mayoritas 3 model.

\section*{5. Clustering}

\subsection*{5a. Clustering vs Classification}
\soal Jelaskan perbedaan clustering dan classification.

\jawab \textbf{Clustering} adalah unsupervised learning: data tidak memiliki label, lalu sistem mencari kelompok berdasarkan kemiripan. \textbf{Classification} adalah supervised learning: data latih memiliki label, lalu model memprediksi label untuk data baru.

\subsection*{5b. Elbow Method untuk K-Means}
\soal Jelaskan cara menentukan nilai $k$ optimal pada k-means. Berdasarkan Gambar 3, tentukan nilai $k$ optimal.

\textbf{Deskripsi Gambar 3:} grafik elbow method menampilkan Sum of Squared Distances terhadap jumlah cluster $k$. Penurunan sangat tajam terjadi dari $k=2$ ke $k=3$, kemudian kurva mulai melandai.

\jawab Nilai $k$ optimal dipilih pada titik siku/elbow, yaitu saat penambahan cluster mulai memberi penurunan SSE yang jauh lebih kecil. Dari grafik, titik siku paling jelas berada di:
\[
\boxed{k=3}
\]
Jadi jumlah cluster optimal adalah sekitar \textbf{3 cluster}.

\end{document}
